\section{Entrega 1: retomada e atualização}
\subsection{Problema, contexto e justificativa}
\subsubsection{Contexto da situação-problema}
A comunicação entre os membros da comunidade escolar é realizada atualmente por diferentes meios, como grupos de WhatsApp, quadros de avisos físicos distribuídos pela escola e comunicados enviados aos responsáveis. Esses canais funcionam de forma descentralizada, sem uma padronização do fluxo de comunicação. Como resultado, informações importantes nem sempre chegam aos destinatários corretos ou são recebidas em tempo hábil, dificultando a participação dos pais na vida escolar dos alunos e prejudicando a divulgação de informações relevantes pela escola.

\subsubsection{Problema ou oportunidade claramente delimitado}
A escola enfrenta dificuldades para garantir que informações institucionais, pedagógicas e administrativas sejam distribuídas de forma organizada, tempestiva e acessível para todos os membros da comunidade escolar. A ausência de um processo estruturado de comunicação reduz a eficácia dos avisos e dificulta o acompanhamento das atividades escolares por alunos e responsáveis.

\subsubsection{Evidências, manifestações e consequências observadas}
\begin{itemize}
\item Avisos importantes são perdidos em meio ao grande volume de mensagens dos grupos de WhatsApp.
\item Nem todos os membros da comunidade escolar participam dos grupos existentes.
\item Comunicados em quadros de avisos físicos possuem alcance limitado e dependem da presença do destinatário na escola.
\item Responsáveis que trabalham possuem dificuldades para acompanhar as informações escolares.
\item Há incerteza sobre quem recebeu ou tomou conhecimento de determinados avisos.
\item O acompanhamento da vida escolar dos alunos ocorre de maneira fragmentada, exigindo consulta a diferentes canais de informação.
\end{itemize}

\subsubsection{Justificativa}
A comunicação eficiente é essencial para o funcionamento das atividades escolares e para o fortalecimento da relação entre escola, alunos e responsáveis. Problemas de comunicação podem causar perda de prazos, baixa participação em eventos, dificuldade de acompanhamento do desempenho escolar e redução do engajamento das famílias. Melhorar esse processo contribui para aumentar a transparência institucional, fortalecer a participação das famílias e apoiar a gestão educacional.

\subsection{Atualização da Entrega 1}
Não foram observadas mudanças nas características apresentadas na Entrega 1 do Projeto Integrador. Permanecem válidos o problema, os objetivos, a visão sistêmica e sociotécnica, o processo AS-IS, os dados e informações iniciais, o escopo e as questões em aberto.
